% !TEX root = ../../../main.tex
% Appendix entry: preprocessing.decompose -- "Decompose" (Appendix A).
\entryfield{Algorithm}
\code{_run} in \code{Adapters/preprocessing_decompose.py} calls \code{decompose} in
\code{Working/Preprocessing/decompose.py} and passes one of its three arrays on. The core adds
no new filter: it reuses the core's running median (\code{make_smooth} in
\code{Working/Detection/analysis/data_analysis.py}), the Detrend block's rolling mean
(\code{make_detrend_filter} in \code{Working/Preprocessing/detrend.py}) and the Highpass
block's Butterworth (\code{make_highpass_filter} in
\code{Working/Detection/analysis/freq_analysis.py}). For a span $x$ of $n$ samples at $f_s$\,Hz,
window $W$ = \emph{baseline\_window\_s} and order $p$ = \emph{order}:
\begin{enumerate}
\item \emph{Refusals.} Any NaN or infinite sample: refused, with the count. $n \le 3(p+1)$:
  refused (the forward-backward filter's padding needs more). A cut $c \ge f_s/2$: refused.
\item \emph{Cut.} $c$ = \emph{noise\_cut\_hz}, or $c = f_s/8$ (a quarter of Nyquist,
  \code{auto_noise_cut_hz}) when it is 0.
\item \emph{Baseline} $b$. With \texttt{rolling\_median} (default): half-width
  $j = \max(1, \operatorname{round}(W f_s/2))$ (\code{baseline_half_width}), window $2j+1$
  samples. If $n > j$, the span is reflect-padded by $j$ samples each side
  (\code{np.pad(..., mode="reflect")}), passed through \code{make_smooth(j)}
  (\code{scipy.ndimage.median_filter}, size $2j+1$, \texttt{mode="nearest"}) and trimmed back to
  $n$:
  \[ b_i = \operatorname{median}\{\tilde x_{i-j},\ldots,\tilde x_{i+j}\} \]
  with $\tilde x$ the reflected span. If $n \le j$, \code{make_smooth(j)} runs on the span
  itself, which pads with copies of the end samples. With \texttt{rolling\_mean}:
  $b = x - \text{detrend}(x)$, the Detrend block's centred rolling mean (\code{_rolling_mean}:
  half-width $\lfloor \operatorname{round}(W f_s)/2\rfloor$, the window shrinking to what fits
  at the ends).
\item \emph{Noise} $\nu = \operatorname{filtfilt}(B, A, x - b)$, with $(B, A)$ a Butterworth
  highpass of order $p$ at $c/(f_s/2)$ (\code{scipy.signal.butter}, \code{filtfilt}: zero
  phase, the effective order doubled by the two passes).
\item \emph{Signal} $s = x - b - \nu$.
\item \emph{Pass on}: the block returns $b$, $\nu$ or $s$ as chosen by \emph{pass\_on}, at
  the same $f_s$.
\end{enumerate}
By construction $b + \nu + s = x$ up to rounding; the block measures
$\max_i |b_i + \nu_i + s_i - x_i|$ and reports it as the reconstruction error.

\paragraph{The half-window rule.} A running median of $2j+1$ samples ignores a departure that
occupies fewer than about half its samples, so an event shorter than about $W/2$ stays in the
signal; one that lasts longer pulls the median towards itself and goes partly to the baseline.
A long, slow recovery counts as part of the event's length. The test
\code{test_an_event_longer_than_about_half_the_window_goes_to_the_baseline} pins this on a
drop of depth 1 on an AR(1) background: with a 60\,s fall and a 600\,s recovery, the signal
layer keeps less than half the depth at its trough at $W = 600$\,s (0.41 measured, per the
test's comment) and more than 0.7 at $W = 3600$\,s; a 10\,s / 60\,s drop keeps more than 0.9
at 600\,s. The help text therefore advises 1800--3600\,s for drops that recover over 5--10
minutes. The median is not a linear filter, so the baseline has no exact frequency edge.

\paragraph{What the block page does with it.} The bridge (\code{_split_layers} in
\code{webui/server/serialize.py}, through the wavelet block's \code{_signal_layers}) ships each
layer as a min/max envelope on the channel's absolute seconds, scaled to the display unit,
with \code{layer_param} (\texttt{pass\_on}) and each row's \code{param_value} (its name). On the
block page, \code{DecompositionLayers} in \code{webui/client/src/analyse/views/registry.tsx}
draws the input and the three layers under \ui{Every layer}, fastest first (noise, signal,
baseline), each on its own $y$ with a scale bar, the chosen one highlighted, with a ``How the
layers were made'' info tip and a note line (\code{SignalSplitNote} in
\code{webui/client/src/analyse/views/SignalView.tsx}: the layer that goes on, the window in
seconds and samples, the cut and whether it was automatic, and the reconstruction error). A
click on a row calls the same \code{setParam} as the Parameters panel with that row's
\code{param_value}: \emph{pass\_on} changes, the row shows a pending outline, and the step is
marked stale (\code{markStale}); it is the next run (\ui{Re-run from}) that sends the new layer
on. Only a payload that names \code{layer_param} gets clickable rows; the wavelet block's do
not.

\entryfield{Assumptions}
\begin{itemize}
\item \emph{Units.} The block is not unit-aware: \code{_run} takes no \code{units} argument and
  every layer is in the recording's stored unit. No setting is an amplitude (the window is in
  seconds, the cut in hertz), so the split does not depend on the unit. The reconstruction
  error in the note is in the stored unit; the drawn layers are in the display unit.
\item \emph{Sampling rate.} The window becomes samples through $f_s$ and the cut is a fraction
  of Nyquist, so both mean the same in seconds and hertz at any rate. The automatic cut moves
  with the rate: 0.125\,Hz at 1\,Hz, 1.25\,Hz at 10\,Hz.
\item \emph{Missing values.} None allowed: a span holding a NaN or inf is refused with the
  count (``Choose a span without gaps''), since one missing sample would spread through the
  median's neighbours and the whole filter output. Nothing is filled.
\item \emph{Span length.} More than $3(p+1)$ samples (16 at the default order 4). A span no
  longer than the half-window still decomposes, with end-sample padding instead of reflection
  (\code{test_a_span_shorter_than_half_the_window_still_decomposes}); its baseline is then
  close to the span's own median.
\item \emph{Signal shape.} Events shorter than about $W/2$, sitting on a wander slower than
  $W$, with fuzz faster than $c$. Where these scales overlap, the layers share the event.
\end{itemize}

\entryfield{Outputs}
A \texttt{Signal} of length $n$ at $f_s$: the chosen layer. \code{meta} holds \code{layers}
(three rows in drawing order noise, signal, baseline, each with \code{name}, \code{kind},
\code{label}, \code{low_hz}/\code{high_hz} where defined, \code{chosen}, \code{rms},
\code{param_value} and an \code{envelope} of at most $2\times1200$ points,
\code{ENVELOPE_BUCKETS}, every sample when it fits), \code{layer} and \code{layer_label} (the
chosen one), \code{layer_param} (\texttt{"pass\_on"}), \code{reconstruction_error}, and
\code{split}: the method, \emph{pass\_on}, the estimator, the window in seconds and in samples,
the cut and \code{noise_cut_auto}, the order, an \code{edge_note} and a \code{nonlinear_note}
(the half-window rule in words). Nothing is written to the database or to disk by the block;
the step cache keeps the output like any other step's.

\entryfield{Cost}
One running median over the span (window $2j+1$, 601 samples at the defaults on 1\,Hz data) and
one forward-backward IIR filter, which is linear in $n$; the median's cost also grows with the
window. Memory is a few arrays of $n$ floats plus the padded copy. The block declares no cost
estimate and no maximum span, and no measurement of its run time is recorded in the
repository.

\entryfield{Failure modes}
\begin{itemize}
\item \emph{A slow event goes to the baseline.} With the window too short for the event (the
  half-window rule), the signal layer keeps only part of it. Check the baseline row.
\item \emph{A sharp edge goes partly to the noise.} The noise layer is a linear highpass of
  $x - b$; a step or a fast fall has content above any cut, so part of the edge appears in the
  noise and the signal's edge is softened. The synthetic test checks only that 5-minute
  Hanning-shaped drops do not leak (correlation of noise with the drops below 0.1).
\item \emph{Edges.} With the median, the span's ends are read reflected: an event the end cuts
  off stays in the signal (\code{test_a_drop_cut_off_by_the_span_end_stays_in_the_signal}), but a
  steady trend running into an end is followed less closely there; the commit records a 1-per-hour
  ramp's baseline lagging by 0.042 at the end. With the rolling mean, the window shrinks at the
  ends. \texttt{filtfilt} adds a short transient at both ends. The \code{edge_note} says to read
  the first and last half-window with care.
\item \emph{Rolling mean baseline.} \texttt{rolling\_mean} is pulled by any event, short or long,
  in proportion to its area; it is offered so a chain that used the Detrend block can see what
  that block removed, not as the better estimator.
\item \emph{Refusals}: a gap, a span of $3(p+1)$ samples or fewer, a cut at or above Nyquist.
\end{itemize}

\entryfield{Validation}
The tests use synthetic signals only. They check that the three layers add up to the input to
$10^{-9}$ on white noise, a random walk, AR(1) with $\phi = 0.98$, a slow drift with drops and a
volts-sized offset, at 1 and 10\,Hz, with both estimators; that each layer is the existing
filter it claims to be; and that on a slow sine plus white fuzz plus 5-minute drops the
baseline follows the sine (RMS error below 0.2), the drops stay in the signal (correlation
above 0.9) and the fuzz goes to the noise (correlation above 0.6). The half-window rule and the
edge handling are pinned as above. No test or evidence script in the repository checks the
split on a real recording.

\entryfield{References}
The method is the project's own arrangement of existing filters (module docstring of
\code{Working/Preprocessing/decompose.py}): a running median for the baseline, and the
Butterworth highpass \citep{butterworth1930theory} as implemented in SciPy
\citep{virtanen2020scipy} for the noise. For a split by octave rather than by role, see
\blockref{preprocessing.wavelet_bands}.

\entryfield{Tests}
\code{tests/test_adapter_decompose.py} (registration and defaults, the sum to the input, each
layer's filter, the automatic cut, the synthetic layer contents, \emph{pass\_on} through a real
chain into Detrend and Invert, the \code{meta} shape, refusals, the half-window rule, the
reflected end and the short span); \code{tests/test_webui_decompose_payload.py} (the layers ship
on absolute seconds in millivolts with the click parameter, survive a step-cache hit, a click
sets the block's parameter, only a payload naming \code{layer_param} gets clickable rows);
\code{tests/test_adapter_spec.py} and \code{tests/test_end_to_end.py} list the block among the
registered adapters.
